% ===========================================================================
\subsection{\file{new-design-reference/css/site.css}}
% ===========================================================================

\textbf{CSS, 768 lines.} Every visual decision on the page: the colours, the
typefaces, the spacing, the layout, the movement, and how all of it rearranges on
a narrow screen. It is the longest file in the repository and, in the judgement
of this document, the most valuable thing in it.

\begin{provenance}{new-design-reference/css/site.css}
Hand-written, and it says where its ideas come from in its own second line:
\code{Design tokens mirror the Digitalise Agency system (tokens.css in Gantry).}
That is a direct statement by the author that the colour and spacing names were
taken from another project in this workspace, \code{gantry-digitalise}, which is
strong evidence of a deliberate shared design language rather than a coincidence.

The file arrived complete in the first commit, \code{0d77d51}, and has not been
edited since except by being moved. No generator banner, no vendor prefix block
of the sort a preprocessor emits, and no minified section: it is written to be
read. \textbf{Confidence: certain for hand-written, high for the token lineage,
which rests on the author's own comment rather than on a file comparison.}
\end{provenance}

\subsubsection*{Why it exists}

Without this file the page still contains every word it needs, and looks like a
1994 document: black serif text on white, everything in one column, no colour, no
layout. Nothing would be broken and nothing would be persuasive. All eight hero
tiles would appear stacked in a list, since it is CSS and not HTML that scatters
them.

\subsubsection*{Line by line}

\paragraph{Lines 1 to 25: the header, and the token block.}

\begin{lstlisting}[firstnumber=1,caption={\file{css/site.css}, lines 1--25}]
/* ai.digitalise.agency · site.css
   Design tokens mirror the Digitalise Agency system (tokens.css in Gantry). */

/* ─── TOKENS ─── */
:root {
  --s0:    #05070d;
  --s1:    #0d1220;
  --s2:    #151c2e;
  --ink1:  #f4f7fc;
  --ink2:  #c3ccdd;
  --ink3:  #93a0b8;
  --accent:  #1967ff;
  --accentH: #3f82ff;
  --accent2: #ff6719;
  --green:   #34d399;
  --border:  #222b40;
  --borderS: #1a2233;
  --font-d: 'Space Grotesk', ui-sans-serif, system-ui, sans-serif;
  --font-b: 'Inter', ui-sans-serif, system-ui, sans-serif;
  --r:   16px;
  --rsm: 10px;
  --ease: cubic-bezier(.22,1,.36,1);
  --maxw: 1320px;
  --shadow: 0 8px 40px rgba(0,0,0,.45);
}
\end{lstlisting}

Lines 1 and 2 are a comment. In CSS, anything between \code{/*} and \code{*/} is
ignored, and unlike HTML there is only this one comment form. These two lines
carry the file's entire provenance, which is why they matter out of all
proportion to their length.

Line 4 is the first of eleven section banners drawn with box-drawing characters.
As with the HTML, they are the author's navigation aid.

Line 5 opens the token block. \code{:root} is a selector that matches the
document's outermost element, which is \code{<html>}. Custom properties defined
there are inherited by every element on the page, which is precisely why token
blocks are written against \code{:root}: it is the one place a definition reaches
everything.

The twenty tokens fall into six groups, and understanding the naming scheme is
most of understanding the file:

\begin{table}[H]\centering\small
\begin{tabular}{@{}llp{6.4cm}@{}}
\toprule
Tokens & Line & What the group is for \\
\midrule
\code{--s0} to \code{--s2}      & 6--8   & Surfaces, darkest to lightest. \code{s0} is the page, \code{s1} is a card, \code{s2} is a detail inside a card \\
\code{--ink1} to \code{--ink3}  & 9--11  & Text, brightest to dimmest. \code{ink1} is a heading, \code{ink2} is body text, \code{ink3} is a caption \\
\code{--accent}, \code{--accentH}, \code{--accent2} & 12--14 & The brand blue, its hover variant, and an orange \\
\code{--green}                  & 15     & The one success colour \\
\code{--border}, \code{--borderS} & 16--17 & A visible hairline, and a subtler one \\
\code{--font-d}, \code{--font-b} & 18--19 & The display typeface and the body typeface \\
\code{--r}, \code{--rsm}, \code{--ease}, \code{--maxw}, \code{--shadow} & 20--24 & Corner radius, small radius, motion curve, content width, drop shadow \\
\bottomrule
\end{tabular}
\caption{The token vocabulary. The numeric suffix always means the same thing:
0 is furthest back or brightest, and the number rises as you come forward or
dim down.}
\end{table}

\begin{keyidea}{Why a numbered scale beats named colours}
Calling a colour \code{--dark-navy} ties the name to how it looks. Calling it
\code{--s1} ties the name to what it is \emph{for}: the surface one step up from
the page. Redesign the site in a light theme and \code{--s1} still means the same
thing while \code{--dark-navy} becomes a lie. This is the single most transferable
idea in this repository, and it is why the author borrowed the scheme from
another project rather than inventing one.
\end{keyidea}

Line 18 defines the display typeface as Space Grotesk with three fallbacks after
it. A browser reads that list left to right and uses the first one it can find.
\code{ui-sans-serif} means the interface font of the operating system,
\code{system-ui} is the same idea by an older name, and \code{sans-serif} is the
last resort every browser guarantees. This is the chain that keeps the page
readable when the Google Fonts request fails.

Line 22 defines the motion curve. \code{cubic-bezier(.22,1,.36,1)} describes how
an animation's speed varies across its duration, by giving two control points of a
curve. This particular one starts fast and settles gently, the shape usually
described as ``ease out''. Every transition in the file refers to it, so all
movement on the page shares one personality.

\begin{honest}{Two tokens are defined and never used}
\code{--accent2} (line 14, an orange) and \code{--shadow} (line 24) are read zero
times in the entire file. Counted with \code{grep -c 'var(--accent2)'}, which
returns \code{0}, and likewise for \code{--shadow}.

The orange is not entirely absent from the page: it appears as the literal
\code{\#ff6719} in the second avatar gradient (line 320) and \code{\#ff9a5c} in
the second proof avatar (line 580). So the colour is used and the token for it is
not, which is the worst of both worlds. Either delete the token or use it.
\end{honest}

\paragraph{Lines 26 to 43: the reset, and the one utility class.}

\begin{lstlisting}[firstnumber=26,caption={\file{css/site.css}, lines 26--43}]

*, *::before, *::after { box-sizing: border-box; margin: 0; padding: 0; }
html { color-scheme: dark; scroll-behavior: smooth; }
body {
  font-family: var(--font-b);
  background: var(--s0);
  color: var(--ink1);
  -webkit-font-smoothing: antialiased;
  text-rendering: optimizeLegibility;
  overflow-x: hidden;
  line-height: 1.5;
}
img, video, svg { max-width: 100%; display: block; }
a { color: inherit; text-decoration: none; }
button { cursor: pointer; border: none; font-family: inherit; }

/* ─── UTILITY ─── */
.container { max-width: var(--maxw); margin: 0 auto; padding: 0 24px; }
\end{lstlisting}

Line 27 is a reset, and it is doing three separate jobs. The selector \code{*}
means every element; \code{*::before} and \code{*::after} extend that to the two
invisible pseudo-elements CSS lets you attach to anything.

\code{box-sizing: border-box} is the most consequential declaration in the file.
By default, if you give an element a width of 200 pixels and then 20 pixels of
padding, it occupies 240 pixels: the padding is added outside the stated width.
That default surprises everyone. \code{border-box} changes it so the stated width
is the final width and padding eats into it. Every layout calculation in this
file assumes it.

\code{margin: 0; padding: 0} removes the browser's built-in spacing on headings,
paragraphs and lists, so that all spacing in the design is deliberate rather than
inherited from a default that varies between browsers.

Line 28 does two things to the document root. \code{color-scheme: dark} tells the
browser this page is dark, so that scrollbars, text-selection highlights and form
controls are drawn in their dark variants rather than staying stubbornly white.
\code{scroll-behavior: smooth} makes every jump-link glide rather than snap.

\begin{honest}{The smooth scroll is implemented twice}
Line 28 already makes anchor links glide. \file{main.js} lines 70 to 77 then
intercept every anchor link and call \code{scrollIntoView(\{behavior: 'smooth'\})}
to do the same thing again. Either alone is sufficient. The JavaScript version is
the redundant one, and it is also the one that carries a latent bug, described in
section 7.3.
\end{honest}

Lines 29 to 37 set the page's defaults: the body typeface, the darkest surface as
the background, the brightest ink as the text, and a line height of 1.5, meaning
each line of text occupies one and a half times the font's size. That last number
is the most important typographic decision in the file and the one most often set
badly.

\code{-webkit-font-smoothing: antialiased} on line 33 thins text rendering on Mac
displays. It is a Safari and Chrome specific property with no standard equivalent,
and it is the conventional companion to light text on a dark background, which
otherwise looks heavier than intended.

\code{overflow-x: hidden} on line 35 forbids horizontal scrolling. Here it is
load-bearing rather than cosmetic: the hero tiles are positioned by percentage and
some of them extend past the right edge of the window, which would otherwise let
the visitor scroll sideways into empty space.

\begin{honest}{Hiding horizontal overflow treats a symptom}
\code{overflow-x: hidden} on the body is a common fix and a blunt one: it hides
the consequence of something overflowing rather than stopping the overflow. The
cost is that if a genuine layout bug ever pushes content off screen, it will be
invisible instead of obvious. Here the cause is known and intentional (the drifting
tiles), so the trade is defensible, but it is worth recognising it as a patch.
\end{honest}

Line 38 makes images, videos and SVGs never exceed their container's width, and
makes them block-level. That second part fixes a genuinely confusing default:
images are inline by default, which reserves space beneath them for the descenders
of text that is not there, producing a mysterious few pixels of gap under every
image.

Line 39 strips links of their blue colour and underline, making them inherit
whatever colour surrounds them. That is a deliberate design choice and an
accessibility trade: colour alone now distinguishes a link from text, which fails
the usual guidance. In this file every link is either a button or a footer item, so
the practical harm is small.

Line 40 styles buttons, and is dead in the strict sense that this page contains no
\code{<button>} element at all. It costs nothing and would come alive the moment
one was added.

Line 43 is the only utility class in the file: a centred column no wider than 1320
pixels with 24 pixels of breathing room on each side. \code{margin: 0 auto} is the
standard centring idiom, meaning no vertical margin and equal automatic margins
left and right. Every section on the page wraps its content in this.

\paragraph{Lines 45 to 75: the button system.}

\begin{lstlisting}[firstnumber=44,caption={\file{css/site.css}, lines 44--75}]

.btn {
  display: inline-flex;
  align-items: center;
  gap: 8px;
  padding: 13px 28px;
  border-radius: 100px;
  font-family: var(--font-d);
  font-size: 15px;
  font-weight: 600;
  letter-spacing: -0.01em;
  transition: transform 0.18s var(--ease), box-shadow 0.18s var(--ease),
              background 0.18s var(--ease);
}
.btn-primary {
  background: var(--accent);
  color: #fff;
}
.btn-primary:hover {
  background: var(--accentH);
  transform: translateY(-2px);
  box-shadow: 0 8px 24px rgba(25,103,255,.4);
}
.btn-ghost {
  background: transparent;
  color: var(--ink2);
  border: 1px solid var(--border);
}
.btn-ghost:hover {
  background: var(--s1);
  color: var(--ink1);
}
\end{lstlisting}

This is the two-class pattern the HTML uses everywhere: \code{.btn} holds
everything both buttons share, and \code{.btn-primary} or \code{.btn-ghost} adds
only what differs. Add a third button style and you write four lines, not thirty.

\code{display: inline-flex} on line 46 makes the button a flex container that sits
in the flow of text rather than on its own line. That matters because the primary
buttons contain both a word and an icon, and flex is what lines those up.

\begin{jargon}{Flexbox}
A layout system that arranges an element's children in a row or a column and
controls how they align and how leftover space is shared. Set
\code{display: flex} on a parent and its children become flex items.
\code{align-items: center} lines them up on their centres, which is how the icon
in a button sits level with the text rather than on the baseline.
\end{jargon}

\code{gap: 8px} on line 48 puts eight pixels between the word and the icon
without either of them needing a margin. Before \code{gap} existed this required
fiddly rules targeting all-but-the-first child, and its arrival is one of the
genuine improvements in modern CSS.

\code{border-radius: 100px} on line 50 makes a pill. Any radius larger than half
the element's height produces a perfect semicircular end, so 100 is simply a
number comfortably larger than needed.

Lines 55 and 56 declare the transition: three named properties, each over 0.18
seconds, each using the shared curve. Listing properties explicitly rather than
writing \code{transition: all} is the better practice, because \code{all} makes
the browser watch every property for changes and can animate things you never
intended.

Lines 62 to 66 are the hover state. \code{transform: translateY(-2px)} lifts the
button two pixels, and the shadow beneath it grows: together they read as the
button rising towards the cursor. The shadow colour \code{rgba(25,103,255,.4)} is
the brand blue at 40\% opacity, written as red, green and blue numbers plus an
alpha.

\begin{jargon}{rgba and opacity}
\code{rgba(25,103,255,.4)} names a colour by its red, green and blue amounts from
0 to 255, plus an opacity from 0 (invisible) to 1 (solid). It is the same blue as
\code{\#1967ff}, expressed in the form that allows transparency.
\end{jargon}

\begin{honest}{Buttons have hover states and no focus states}
Neither button defines \code{:focus} or \code{:focus-visible}. A visitor
navigating with a keyboard rather than a mouse therefore gets whatever outline the
browser draws by default, which on a dark background is often nearly invisible.
Since these are the page's only interactive elements, this is the most consequential
accessibility gap in the stylesheet, and it is three lines to fix.
\end{honest}

\paragraph{Lines 77 to 92: the two reveal mechanisms.}

\begin{lstlisting}[firstnumber=76,caption={\file{css/site.css}, lines 76--92}]

/* ─── REVEAL (scroll-driven, driven by JS adding .visible) ─── */
.reveal {
  opacity: 0;
  transform: translateY(28px);
  transition: opacity 0.7s var(--ease), transform 0.7s var(--ease);
}
.reveal.visible { opacity: 1; transform: none; }

/* stagger via CSS custom property set inline */
.stagger > * {
  opacity: 0;
  transform: translateY(20px);
  transition: opacity 0.6s var(--ease), transform 0.6s var(--ease);
  transition-delay: calc(var(--i, 0) * 90ms);
}
.stagger.visible > * { opacity: 1; transform: none; }
\end{lstlisting}

These sixteen lines are the entire animation contract between the stylesheet and
the JavaScript, and the comment on line 77 states it: the CSS defines both states
and the transition between them, and the JavaScript's only job is to add the class
\code{visible} at the right moment.

\code{.reveal} starts invisible and 28 pixels low. \code{.reveal.visible}, which
matches an element carrying \emph{both} classes, returns it to fully visible and
its natural position. Because a transition is declared, the browser animates
between the two over 0.7 seconds rather than snapping.

The stagger is the same idea with one addition. \code{.stagger > *} means every
direct child of a staggered container, and line 90 delays each child's transition
by \code{calc(var(--i, 0) * 90ms)}: the number the HTML wrote inline, multiplied by
90 milliseconds. Child 0 starts immediately, child 3 starts 270 milliseconds later.

\begin{jargon}{calc, and a custom property's fallback}
\code{calc()} performs arithmetic in CSS, and can mix units.
\code{var(--i, 0)} reads the custom property \code{--i} and falls back to
\code{0} if it was never set, which is what keeps this rule safe on any element
whose inline \code{--i} was forgotten: it simply gets no delay instead of breaking.
\end{jargon}

\begin{keyidea}{Why the animation lives in CSS and the trigger lives in JavaScript}
This split is the reason the page degrades well. If JavaScript never runs, the
class \code{visible} is never added, and every revealed element stays at opacity
zero, invisible forever. That would be catastrophic, and it is exactly why line
767 exists: the reduced-motion block at the end of the file forces
\code{.reveal} and \code{.stagger > *} back to full opacity. That block is doing
double duty as an accessibility feature and as a safety net.
\end{keyidea}

\paragraph{Lines 94 to 133: the navigation bar.}

\begin{lstlisting}[firstnumber=93,caption={\file{css/site.css}, lines 93--133}]

/* ─── NAV ─── */
.nav {
  position: fixed;
  top: 0; left: 0; right: 0;
  z-index: 100;
  padding: 18px 24px;
  display: flex;
  align-items: center;
  justify-content: space-between;
  transition: background 0.3s, border-color 0.3s, box-shadow 0.3s;
}
.nav.scrolled {
  background: rgba(5,7,13,0.88);
  border-bottom: 1px solid var(--border);
  box-shadow: 0 4px 24px rgba(0,0,0,.3);
  backdrop-filter: blur(16px);
  -webkit-backdrop-filter: blur(16px);
}
.nav-inner {
  max-width: var(--maxw);
  width: 100%;
  margin: 0 auto;
  display: flex;
  align-items: center;
  justify-content: space-between;
}
.nav-logo {
  display: flex;
  align-items: center;
  gap: 10px;
  font-family: var(--font-d);
  font-size: 18px;
  font-weight: 700;
  letter-spacing: -0.02em;
  color: var(--ink1);
}
.nav-logo .logo-ai {
  color: var(--accent);
}
.nav-logo svg { width: 28px; height: 28px; }
\end{lstlisting}

\code{position: fixed} on line 96 pins the bar to the window rather than to the
document, so it stays put while the page scrolls beneath it. Lines 97 pins it to
the top and stretches it across.

\code{z-index: 100} on line 98 controls stacking order: when elements overlap, the
higher number is drawn in front. 100 is arbitrary and generous, chosen so nothing
else on the page can accidentally cover the bar.

Lines 105 to 111 are the scrolled state, the class the JavaScript adds once the
page has moved more than 40 pixels. It gives the bar a background at 88\% opacity,
a bottom hairline, and a shadow.

\code{backdrop-filter: blur(16px)} on line 109 is the interesting one: it blurs
whatever is \emph{behind} the element, producing the frosted-glass effect that only
works because the background is slightly transparent. Line 110 repeats it with a
\code{-webkit-} prefix for Safari, which needed the prefixed form for years.

\begin{honest}{A doubled flex layout}
\code{.nav} sets \code{display: flex} with \code{justify-content: space-between}
(lines 100 to 102), and \code{.nav-inner} sets exactly the same three properties
again (lines 116 to 118). The outer one has only one child, so its
\code{space-between} has nothing to space. It is harmless and it is redundant, and
it makes the layout harder to reason about than it needs to be.
\end{honest}

Line 130 is a descendant selector: \code{.nav-logo .logo-ai}, with a space between
the two names, means ``an element with class \code{logo-ai} anywhere inside an
element with class \code{nav-logo}''. That space is meaningful syntax, and reading
it as one thing rather than two is a common beginner error.

\paragraph{Lines 135 to 177: the hero shell and the floating animation.}

\begin{lstlisting}[firstnumber=134,caption={\file{css/site.css}, lines 134--177}]

/* ─── HERO ─── */
.hero {
  position: relative;
  min-height: 100vh;
  display: flex;
  align-items: center;
  justify-content: center;
  overflow: hidden;
  padding-top: 80px;
}

/* animated tile backdrop */
.hero-bg {
  position: absolute;
  inset: 0;
  pointer-events: none;
}
.tile {
  position: absolute;
  width: 176px;
  background: var(--s1);
  border: 1px solid var(--border);
  border-radius: var(--r);
  padding: 14px;
  opacity: 0;
  animation: tileFloat 12s var(--ease) infinite;
}
.tile:nth-child(1) { top: 12%; left: 3%;   animation-delay:  0.0s; }
.tile:nth-child(2) { top: 18%; right: 4%;  animation-delay:  1.5s; }
.tile:nth-child(3) { top: 55%; left: 5%;   animation-delay:  3.0s; }
.tile:nth-child(4) { bottom: 18%; right: 3%; animation-delay: 4.5s; }
.tile:nth-child(5) { top: 42%; right: 26%; animation-delay:  6.0s; }
.tile:nth-child(6) { top:  8%; left: 32%;  animation-delay:  7.5s; }
.tile:nth-child(7) { bottom: 28%; left: 22%; animation-delay: 9.0s; }
.tile:nth-child(8) { top: 30%; left: 64%;  animation-delay: 10.5s; }

@keyframes tileFloat {
  0%   { opacity: 0; transform: translateY(18px) scale(0.94); }
  12%  { opacity: 0.8; transform: translateY(0) scale(1); }
  72%  { opacity: 0.8; transform: translateY(-10px) scale(1); }
  85%  { opacity: 0; transform: translateY(-18px) scale(0.94); }
  100% { opacity: 0; transform: translateY(-18px) scale(0.94); }
}
\end{lstlisting}

\code{position: relative} on line 137 looks like it does nothing, and it is
essential. It establishes this element as the reference point for any absolutely
positioned descendant. Without it, the tiles would position themselves against the
whole window instead of against the hero.

\begin{jargon}{Positioning: relative and absolute}
\code{position: absolute} takes an element out of the normal flow and places it at
coordinates you give. Those coordinates are measured from the nearest ancestor
that is itself positioned, which in practice means the nearest one with
\code{position: relative}. This pairing, a relative parent and absolute children,
is the standard way to place things inside a box.
\end{jargon}

\code{min-height: 100vh} on line 138 makes the hero at least as tall as the
window. \code{vh} means one percent of the viewport height, so \code{100vh} is the
full height.

\code{pointer-events: none} on line 149 makes the tile container transparent to the
mouse: clicks pass straight through to whatever is beneath. Without it, the tiles
would sit over the hero buttons and swallow clicks aimed at ``Book a call''. It is
a one-line fix for a bug that would otherwise be maddening to diagnose, since the
button would look perfectly normal and simply not respond.

Lines 152 to 161 define a tile: absolutely positioned, 176 pixels wide, using the
card surface, the standard border and the standard 16-pixel radius, and starting
at \code{opacity: 0}. It is invisible until the animation makes it visible.

Line 160 runs the animation: the named sequence \code{tileFloat}, over 12 seconds,
using the shared curve, repeating forever.

Lines 162 to 169 place the eight tiles individually. \code{:nth-child(1)} matches
the first child of its parent, and so on.

\begin{keyidea}{This is where the HTML order becomes load-bearing}
Nothing here mentions \code{tile-blue} or \code{tile-orange}. The tiles are
positioned purely by their \emph{number}, so tile 5 is wherever the fifth child
happens to be. Reorder the blocks in \file{index.html} and the layout of the hero
changes, with no change to this file. That coupling is invisible from either file
alone, and it is the kind of thing this document exists to make visible.
\end{keyidea}

Note also that the eight delays are 1.5 seconds apart, from 0.0 up to 10.5, across
a 12-second loop. That is a deliberate arrangement: eight evenly spaced entrances
inside one cycle, so at any moment several tiles are visible at different stages
and the backdrop never appears empty or synchronised.

Lines 171 to 177 define the animation itself.

\begin{jargon}{Keyframes}
An animation defined by describing what the element looks like at particular
points through its run, given as percentages of the total duration. The browser
fills in every moment in between. \code{0\%} is the start and \code{100\%} is the
end.
\end{jargon}

Read as a story: at the start the tile is invisible, 18 pixels low and slightly
small. By 12\% of the way through (about 1.4 seconds) it has risen into place at
full size and 80\% opacity. It holds there until 72\%, drifting up 10 pixels. By
85\% it has faded out completely while continuing to rise and shrink. The final
15\% is spent invisible, which is the gap before it begins again.

\paragraph{Lines 179 to 218: what is inside a tile.}

\begin{lstlisting}[firstnumber=178,caption={\file{css/site.css}, lines 178--218}]

.tile-thumb {
  border-radius: 10px;
  height: 90px;
  margin-bottom: 10px;
  display: flex;
  align-items: center;
  justify-content: center;
  position: relative;
  overflow: hidden;
}
.tile-thumb svg { width: 28px; height: 28px; color: rgba(255,255,255,.6); position: relative; z-index: 1; }
.tile-meta { display: flex; align-items: center; justify-content: space-between; gap: 6px; }
.tile-label { font-size: 11px; color: var(--ink3); font-weight: 500; }
.tile-metric {
  font-family: var(--font-d);
  font-size: 11px;
  font-weight: 700;
  color: var(--green);
  background: rgba(52,211,153,.1);
  border-radius: 20px;
  padding: 2px 7px;
}

/* tile gradient variants */
.tile-blue .tile-thumb  { background: linear-gradient(135deg,#0a1e50,#0d2c6e); }
.tile-orange .tile-thumb { background: linear-gradient(135deg,#3a1800,#5a2800); }
.tile-teal .tile-thumb  { background: linear-gradient(135deg,#003838,#005252); }
.tile-violet .tile-thumb { background: linear-gradient(135deg,#1a0a40,#2a1060); }

/* fake text bars inside "post" tiles */
.tile-lines { padding: 0; }
.tile-line {
  height: 7px;
  border-radius: 4px;
  background: var(--s2);
  margin-bottom: 6px;
}
.tile-line:nth-child(1) { width: 100%; }
.tile-line:nth-child(2) { width: 85%; }
.tile-line:nth-child(3) { width: 60%; }
\end{lstlisting}

Lines 179 to 188 shape the thumbnail area: a fixed 90-pixel-tall rounded box that
centres its contents both ways using flex.

Line 189 styles the icon inside it. Note \code{color:} rather than \code{fill:} or
\code{stroke:}: this works because the SVG in the HTML was written with
\code{stroke="currentColor"}, so setting a text colour here reaches into the
drawing. That is the pay-off for the attribute noted in section 7.1.

Line 190 lays out the caption row with \code{space-between}, which pushes the label
to the left edge and the metric to the right, the flex idiom for a two-ended row.

Lines 192 to 200 style the metric as a small green pill: green text on a green
background at 10\% opacity, which is the standard recipe for a tinted badge that
suits any surface beneath it.

Lines 203 to 206 are the four colour variants. Each is a descendant selector: the
gradient applies to a \code{.tile-thumb} \emph{inside} a \code{.tile-blue}, not to
the tile itself. \code{linear-gradient(135deg, A, B)} blends from colour A to
colour B along a line at 135 degrees, which runs from top-left to bottom-right.

Line 209 is the dead rule identified in section 7.1: both places that use
\code{.tile-lines} override its padding inline, so this never applies.

Lines 210 to 218 draw the fake text bars: 7 pixels tall, rounded, in the \code{s2}
surface, with the three widths that make them read as a paragraph. Note that tile
8 overrides the third width inline to 50\%, so these are defaults rather than
rules.

\paragraph{Lines 220 to 243: the statistic tiles and the readability overlay.}

\begin{lstlisting}[firstnumber=219,caption={\file{css/site.css}, lines 219--243}]

/* stats tile */
.tile-stat-num {
  font-family: var(--font-d);
  font-size: 26px;
  font-weight: 700;
  color: var(--ink1);
  line-height: 1;
  margin-bottom: 4px;
}
.tile-stat-lbl { font-size: 11px; color: var(--ink3); }

/* overlay - keeps center readable */
.hero-overlay {
  position: absolute;
  inset: 0;
  background:
    radial-gradient(ellipse 55% 65% at 50% 50%,
      rgba(5,7,13,.9) 0%,
      rgba(5,7,13,.55) 70%,
      transparent 100%),
    linear-gradient(to bottom,
      rgba(5,7,13,.2) 0%,
      rgba(5,7,13,.85) 100%);
}
\end{lstlisting}

\code{line-height: 1} on line 226 is deliberate. The page-wide default of 1.5 is
right for paragraphs and wrong for a large number, where it leaves a visible gap
above and below. Setting it to 1 makes the line exactly as tall as the digits.

Lines 231 to 243 are the overlay, and they are the most sophisticated piece of
styling in the file. Two gradients are stacked in one declaration, separated by a
comma, and in CSS the \emph{first} one listed is drawn on top.

The first is a radial gradient: an ellipse 55\% wide and 65\% tall, centred, going
from 90\% opaque near-black at its centre, through 55\% at seven-tenths of the way
out, to fully transparent at the edge. That is a dark cloud over the middle of the
hero, exactly where the headline sits.

The second is a linear gradient running down the whole area, from 20\% opaque at
the top to 85\% at the bottom, which darkens the lower half so the hero blends into
the section beneath it.

Together they solve the real problem with a busy animated backdrop: text over
moving content is hard to read, and the contrast changes as the content moves.
Rather than dimming the tiles themselves, the design puts a shaped veil between
them and the text.

\begin{jargon}{inset}
\code{inset: 0} is shorthand for \code{top: 0; right: 0; bottom: 0; left: 0},
which on an absolutely positioned element means ``stretch to fill the parent
completely''.
\end{jargon}

\paragraph{Lines 245 to 324: the hero's content.}

\begin{lstlisting}[firstnumber=244,caption={\file{css/site.css}, lines 244--275}]

.hero-content {
  position: relative;
  z-index: 2;
  text-align: center;
  max-width: 700px;
  padding: 0 24px;
}
.hero-badge {
  display: inline-flex;
  align-items: center;
  gap: 8px;
  background: rgba(25,103,255,.12);
  border: 1px solid rgba(25,103,255,.3);
  border-radius: 100px;
  padding: 6px 14px;
  font-size: 13px;
  font-weight: 600;
  color: var(--accentH);
  letter-spacing: 0.01em;
  margin-bottom: 28px;
}
.hero-badge-dot {
  width: 7px; height: 7px;
  border-radius: 50%;
  background: var(--accent);
  animation: pulse 2s ease-in-out infinite;
}
@keyframes pulse {
  0%, 100% { opacity: 1; transform: scale(1); }
  50%       { opacity: 0.5; transform: scale(0.7); }
}
\end{lstlisting}

\code{z-index: 2} on line 247 lifts the content above the tiles and the overlay,
which is what makes the headline readable rather than buried. Note that
\code{z-index} only works on positioned elements, which is why line 246 sets
\code{position: relative} even though nothing is being moved.

The badge (lines 252 to 265) uses the tinted-pill recipe again: brand blue at 12\%
for the fill and 30\% for the border, with the brighter hover blue for the text so
it stays legible against its own tint.

Lines 266 to 275 define the pulsing dot: a 7-pixel square made circular with a 50\%
radius, running a two-second pulse forever. The keyframes group \code{0\%, 100\%}
on one line, which is the idiom for an animation that returns to where it started:
full size at both ends, shrunk to 70\% and half-faded in the middle.

\begin{lstlisting}[firstnumber=276,caption={\file{css/site.css}, lines 276--324}]
.hero-h1 {
  font-family: var(--font-d);
  font-size: clamp(42px, 6.5vw, 80px);
  font-weight: 700;
  letter-spacing: -0.03em;
  line-height: 1.06;
  color: var(--ink1);
  margin-bottom: 36px;
}
.hero-h1 em {
  font-style: normal;
  color: var(--accent);
}
.hero-actions {
  display: flex;
  align-items: center;
  justify-content: center;
  gap: 14px;
  flex-wrap: wrap;
  margin-bottom: 48px;
}
.hero-social-proof {
  display: flex;
  align-items: center;
  gap: 12px;
  justify-content: center;
}
.avatars { display: flex; }
.avatar {
  width: 34px; height: 34px;
  border-radius: 50%;
  border: 2px solid var(--s0);
  margin-left: -10px;
  display: flex;
  align-items: center;
  justify-content: center;
  font-family: var(--font-d);
  font-size: 12px;
  font-weight: 700;
  color: #fff;
  flex-shrink: 0;
}
.avatars .avatar:first-child { margin-left: 0; }
.av1 { background: linear-gradient(135deg,#1967ff,#3f82ff); }
.av2 { background: linear-gradient(135deg,#ff6719,#ff8c4a); }
.av3 { background: linear-gradient(135deg,#34d399,#10b981); }
.av4 { background: linear-gradient(135deg,#a78bfa,#7c3aed); }
.hero-proof-text { font-size: 13px; color: var(--ink3); }
.hero-proof-text strong { color: var(--ink2); }
\end{lstlisting}

Line 278 is the most useful single line in the file to learn from.

\begin{jargon}{clamp}
\code{clamp(minimum, preferred, maximum)} picks the preferred value, but never lets
it fall below the minimum or rise above the maximum. Here the preferred value is
\code{6.5vw}, meaning 6.5\% of the window width, so the headline scales
continuously with the window between a floor of 42 pixels and a ceiling of 80. One
line replaces what used to need several media queries.
\end{jargon}

Line 280 sets \code{letter-spacing: -0.03em}, pulling the letters slightly closer
together. Large display type looks loose at its default spacing, and negative
tracking at large sizes is standard practice. The unit \code{em} means ``relative
to this element's font size'', so the tightening scales with the text.

Lines 285 to 288 are the rule flagged in section 7.1: it takes the \code{<em>} tag,
removes the italics it normally implies, and colours it blue instead.

Line 294 sets \code{flex-wrap: wrap} on the button row, so on a narrow screen the
second button drops onto its own line instead of overflowing. That single property
is why the hero needs no special mobile rule for its buttons.

Lines 303 to 318 build the overlapping avatar row, and the technique is worth
naming: each circle has a negative left margin of 10 pixels, pulling it over its
neighbour, and a 2-pixel border in the page background colour, which reads as a gap
punched out of the circle behind it. Line 318 then cancels the negative margin on
the first one so the row does not start shifted left.

\code{flex-shrink: 0} on line 316 stops the circles being squashed when space is
tight. Flex items shrink by default, and circles that shrink become ovals.

\begin{honest}{Four literal colours where tokens exist}
Lines 319 to 322 write \code{\#1967ff}, \code{\#3f82ff}, \code{\#ff6719} and
\code{\#34d399} directly. Three of those four have tokens (\code{--accent},
\code{--accentH}, \code{--green}) and the fourth has one that is otherwise unused
(\code{--accent2}). Gradients accept custom properties perfectly well, so there is
no technical reason for the literals. This is the same inconsistency as the inline
styles in the HTML, and together they mean the token system governs most of the
page rather than all of it.
\end{honest}

\paragraph{Lines 326 to 359: the metrics band.}

\begin{lstlisting}[firstnumber=325,caption={\file{css/site.css}, lines 325--359}]

/* ─── METRICS BAND ─── */
.metrics {
  padding: 80px 0;
  border-top: 1px solid var(--borderS);
  border-bottom: 1px solid var(--borderS);
}
.metrics-grid {
  display: grid;
  grid-template-columns: repeat(4, 1fr);
  gap: 2px;
}
.metric-item {
  padding: 36px 24px;
  text-align: center;
  border-right: 1px solid var(--borderS);
}
.metric-item:last-child { border-right: none; }
.metric-num {
  font-family: var(--font-d);
  font-size: clamp(36px, 4.5vw, 60px);
  font-weight: 700;
  letter-spacing: -0.03em;
  line-height: 1;
  color: var(--ink1);
  margin-bottom: 8px;
}
.metric-num .accent { color: var(--accent); }
.metric-label {
  font-size: 13px;
  color: var(--ink3);
  font-weight: 500;
  text-transform: uppercase;
  letter-spacing: 0.06em;
}
\end{lstlisting}

Line 333 introduces the second layout system in the file.

\begin{jargon}{CSS Grid}
A layout system for two dimensions, as opposed to flexbox's one.
\code{grid-template-columns: repeat(4, 1fr)} means four columns, each taking one
equal share of the space. \code{1fr} is a fraction of what is left over after any
fixed sizes are accounted for.
\end{jargon}

The separators here are drawn as a right border on every item (line 340), with the
last one removed (line 342). That is a different technique from the one used later
by the services grid, which draws its lines using gaps. Both work; using both in
one file is an inconsistency, and the border approach is the one that needed a
special case.

Line 352, \code{.metric-num .accent}, styles an element with class \code{accent}
inside a metric number. \textbf{No such element exists} anywhere in
\file{index.html}. This rule is dead.

Lines 353 to 359 style the label beneath each number: small, dim, uppercase, with
positive letter-spacing. Note the tracking is positive here and negative on the
big numbers above. That is not an inconsistency but a rule of typography: large
text needs tightening, small uppercase text needs opening up or it becomes a
smudge.

\paragraph{Lines 361 to 379: shared section headings.}

\begin{lstlisting}[firstnumber=360,caption={\file{css/site.css}, lines 360--379}]

/* ─── PIPELINE ─── */
.pipeline-section { padding: 120px 0; }
.section-kicker {
  font-size: 12px;
  font-weight: 700;
  text-transform: uppercase;
  letter-spacing: 0.1em;
  color: var(--accent);
  margin-bottom: 14px;
}
.section-title {
  font-family: var(--font-d);
  font-size: clamp(30px, 4vw, 52px);
  font-weight: 700;
  letter-spacing: -0.025em;
  line-height: 1.1;
  color: var(--ink1);
  margin-bottom: 72px;
}
\end{lstlisting}

Although these sit under the pipeline banner, \code{.section-kicker} and
\code{.section-title} are used by all four remaining sections. Their placement here
is a small organisational untidiness: a reader looking for the heading styles would
reasonably expect them near the top with the other shared pieces, not inside the
pipeline block.

The kicker is the small blue uppercase line (``How it works'', ``The work'') and the
title is the large heading beneath it. Both follow the same tracking rule as the
metrics band.

\begin{honest}{The kicker's contrast is worth measuring}
\code{--accent} is \code{\#1967ff}, a saturated blue, and the kicker sets it at 12
pixels in uppercase on the near-black page background. Small text needs more
contrast than large text, not less. This document does not have a measured
contrast ratio for that exact pairing and so does not assert a failure, but it is
the first thing to check in an accessibility audit, along with the 11-pixel service
tags which use the dimmest ink on the same background.
\end{honest}

\paragraph{Lines 381 to 450: the pipeline.}

\begin{lstlisting}[firstnumber=380,caption={\file{css/site.css}, lines 380--426}]

.pipeline {
  display: flex;
  align-items: center;
  gap: 0;
  position: relative;
  overflow: hidden;
}
.pipeline-node {
  flex: 1;
  display: flex;
  flex-direction: column;
  align-items: center;
  gap: 16px;
  position: relative;
  z-index: 1;
  opacity: 0;
  transform: scale(0.85);
  transition: opacity 0.5s var(--ease), transform 0.5s var(--ease);
}
.pipeline.visible .pipeline-node:nth-child(1) { opacity: 1; transform: none; transition-delay: 0.0s; }
.pipeline.visible .pipeline-node:nth-child(3) { opacity: 1; transform: none; transition-delay: 0.2s; }
.pipeline.visible .pipeline-node:nth-child(5) { opacity: 1; transform: none; transition-delay: 0.4s; }
.pipeline.visible .pipeline-node:nth-child(7) { opacity: 1; transform: none; transition-delay: 0.6s; }
.pipeline.visible .pipeline-node:nth-child(9) { opacity: 1; transform: none; transition-delay: 0.8s; }

.pipeline-connector {
  flex: 0 0 auto;
  width: 60px;
  height: 2px;
  background: var(--border);
  position: relative;
  overflow: hidden;
}
.pipeline-connector::after {
  content: '';
  position: absolute;
  top: 0; left: 0;
  height: 100%;
  width: 0;
  background: linear-gradient(to right, var(--accent), var(--accentH));
  transition: width 0.6s var(--ease);
}
.pipeline.visible .pipeline-connector:nth-child(2)::after { width: 100%; transition-delay: 0.15s; }
.pipeline.visible .pipeline-connector:nth-child(4)::after { width: 100%; transition-delay: 0.35s; }
.pipeline.visible .pipeline-connector:nth-child(6)::after { width: 100%; transition-delay: 0.55s; }
.pipeline.visible .pipeline-connector:nth-child(8)::after { width: 100%; transition-delay: 0.75s; }
\end{lstlisting}

This is the most intricate animation in the file and it is worth reading closely,
because it is doing something the stagger mechanism could not.

The pipeline is a flex row of nine children alternating node, connector, node,
connector and so on: five nodes at odd positions and four connectors at even ones.
That alternation is why the selectors count 1, 3, 5, 7, 9 for nodes and 2, 4, 6, 8
for connectors.

\code{flex: 1} on the nodes (line 389) tells each to take an equal share of the
leftover space. \code{flex: 0 0 auto} on the connectors (line 407) tells them the
opposite: do not grow, do not shrink, stay exactly 60 pixels. So the nodes spread
and the connectors stay fixed.

Each node begins invisible and at 85\% size, and the \code{.visible} rules bring it
to full size with an increasing delay, 0.2 seconds apart. The connectors do
something different: each has an \code{::after} pseudo-element, an invisible extra
box CSS creates on demand, absolutely positioned, starting at zero width and
filled with a blue gradient. When the section becomes visible its width transitions
to 100\%, so the blue appears to flow along the line.

\begin{jargon}{Pseudo-element}
An extra box that CSS creates inside an element without any corresponding tag in
the HTML, written as \code{::before} or \code{::after}. It requires the
\code{content} property to exist at all, which is why line 415 sets
\code{content: ''}: an empty string is enough to bring it into being. It is the
standard way to draw a decorative shape without polluting the markup.
\end{jargon}

The delays interleave deliberately: node 1 at 0.0, connector at 0.15, node 3 at
0.2, connector at 0.35, and so on. The line starts drawing before the next node
lands, which is what makes the sequence read as flowing rather than as a series of
separate appearances.

\begin{honest}{This animation cannot survive an edit to the HTML}
Every selector here depends on the exact child positions in \file{index.html}. Add
a sixth step, or wrap the nodes in a container, and every one of these nine rules
silently stops matching: the pipeline would simply never appear, staying at opacity
zero forever, with no error anywhere. Selecting by class
(\code{.pipeline-node:nth-of-type(n)}) would have been sturdier for the same
effect.
\end{honest}

\begin{lstlisting}[firstnumber=427,caption={\file{css/site.css}, lines 427--450}]

.pipeline-icon {
  width: 72px; height: 72px;
  border-radius: 20px;
  background: var(--s1);
  border: 1px solid var(--border);
  display: flex;
  align-items: center;
  justify-content: center;
  transition: border-color 0.3s, box-shadow 0.3s;
}
.pipeline-icon svg { width: 30px; height: 30px; color: var(--accent); }
.pipeline-node:hover .pipeline-icon {
  border-color: var(--accent);
  box-shadow: 0 0 24px rgba(25,103,255,.25);
}
.pipeline-lbl {
  font-family: var(--font-d);
  font-size: 13px;
  font-weight: 600;
  color: var(--ink2);
  text-align: center;
  white-space: nowrap;
}
\end{lstlisting}

Lines 439 to 442 give each icon box a blue border and a soft blue glow when the
mouse is over its node. A shadow with no offset and a large blur, as here, reads as
a glow rather than a shadow.

\code{white-space: nowrap} on line 449 stops labels such as ``Auto-Post'' wrapping
onto two lines and disturbing the row's alignment. It is overridden at line 739 for
narrow screens, where the pipeline stacks vertically and there is no alignment left
to protect.

\paragraph{Lines 452 to 539: the portfolio cards.}

\begin{lstlisting}[firstnumber=451,caption={\file{css/site.css}, lines 451--509}]

/* ─── PORTFOLIO ─── */
.portfolio-section { padding: 0 0 120px; }
.portfolio-grid {
  display: grid;
  grid-template-columns: repeat(3, 1fr);
  gap: 16px;
}
.port-card {
  background: var(--s1);
  border: 1px solid var(--border);
  border-radius: var(--r);
  overflow: hidden;
  transition: transform 0.22s var(--ease), border-color 0.22s, box-shadow 0.22s;
  cursor: pointer;
}
.port-card:hover {
  transform: translateY(-4px);
  border-color: rgba(25,103,255,.4);
  box-shadow: 0 16px 48px rgba(0,0,0,.4);
}
.port-thumb {
  aspect-ratio: 16/9;
  position: relative;
  display: flex;
  align-items: center;
  justify-content: center;
  overflow: hidden;
}
.port-thumb-overlay {
  position: absolute;
  inset: 0;
  background: rgba(5,7,13,.25);
  display: flex;
  align-items: center;
  justify-content: center;
  opacity: 0;
  transition: opacity 0.22s;
}
.port-card:hover .port-thumb-overlay { opacity: 1; }
.port-play {
  width: 52px; height: 52px;
  border-radius: 50%;
  background: rgba(255,255,255,.15);
  backdrop-filter: blur(8px);
  display: flex;
  align-items: center;
  justify-content: center;
  border: 1.5px solid rgba(255,255,255,.3);
}
.port-play svg { width: 22px; height: 22px; color: #fff; margin-left: 2px; }

/* gradient thumbnails - replace with real video posters */
.pt1 { background: linear-gradient(135deg,#0a1840 0%,#0d2c6e 50%,#0a1840 100%); }
.pt2 { background: linear-gradient(135deg,#2a0e00 0%,#5a2800 50%,#2a0e00 100%); }
.pt3 { background: linear-gradient(135deg,#001a28 0%,#003855 50%,#001a28 100%); }
.pt4 { background: linear-gradient(135deg,#150050 0%,#2d0080 50%,#150050 100%); }
.pt5 { background: linear-gradient(135deg,#0a2800 0%,#1a4800 50%,#0a2800 100%); }
.pt6 { background: linear-gradient(135deg,#280014 0%,#500028 50%,#280014 100%); }
\end{lstlisting}

\code{overflow: hidden} on line 463 is what makes the rounded corners work: without
it, the gradient thumbnail inside would square off the top two corners it is
supposed to follow.

\code{cursor: pointer} on line 465 changes the mouse to a pointing hand, which
conventionally means clickable. As section 7.1 noted, the card is not clickable,
so this is a promise the page does not keep.

\code{aspect-ratio: 16/9} on line 473 reserves a widescreen-shaped area regardless
of width. This is worth appreciating: it means the layout does not shift when a
real image eventually loads, because the space is already the right shape. It is
the modern replacement for a well-known hack involving percentage padding.

The overlay (lines 480 to 489) sits over the whole thumbnail at zero opacity and
fades to one on hover, revealing the play button beneath a 25\% black wash. The play
button itself (lines 491 to 500) is a frosted circle: white at 15\% opacity with an
8-pixel backdrop blur and a brighter rim.

Line 501 gives the play triangle \code{margin-left: 2px}. That is not an accident.
A triangle's visual centre of mass sits left of its geometric centre, so a
perfectly centred play icon looks slightly left. Nudging it two pixels right is a
standard correction, and noticing it is a good sign that whoever built this cared
about detail.

Lines 504 to 509 are the six placeholder gradients, each a three-stop blend that
returns to its starting colour, so it reads as a soft sheen rather than a slope.
The comment on line 503 says outright that they are standing in for real video
posters that never arrived.

\begin{lstlisting}[firstnumber=510,caption={\file{css/site.css}, lines 510--539}]

.port-badge {
  position: absolute;
  top: 10px; right: 10px;
  background: rgba(5,7,13,.75);
  backdrop-filter: blur(8px);
  border: 1px solid var(--border);
  border-radius: 20px;
  padding: 4px 10px;
  font-size: 11px;
  font-weight: 600;
  color: var(--green);
}
.port-info { padding: 16px; }
.port-platform {
  display: flex;
  align-items: center;
  gap: 6px;
  font-size: 11px;
  color: var(--ink3);
  margin-bottom: 6px;
}
.port-platform svg { width: 13px; height: 13px; }
.port-title {
  font-family: var(--font-d);
  font-size: 15px;
  font-weight: 600;
  color: var(--ink1);
  letter-spacing: -0.01em;
}
\end{lstlisting}

The badge is pinned to the thumbnail's top-right corner, which works because
\code{.port-thumb} set \code{position: relative} on line 474.

The card title on lines 533 to 539 is the only place in the file where the display
typeface is used at a small size (15 pixels). Everywhere else Space Grotesk
appears it is either large (a heading, a metric) or tiny and uppercase (a label),
so this is the one instance of it doing ordinary work.

\paragraph{Lines 541 to 604: the client proof cards.}

\begin{lstlisting}[firstnumber=540,caption={\file{css/site.css}, lines 540--604}]

/* ─── CLIENT PROOF ─── */
.proof-section { padding: 0 0 120px; }
.proof-grid {
  display: grid;
  grid-template-columns: repeat(3, 1fr);
  gap: 16px;
}
.proof-card {
  background: var(--s1);
  border: 1px solid var(--border);
  border-radius: var(--r);
  padding: 28px;
  display: flex;
  flex-direction: column;
  gap: 20px;
  transition: border-color 0.22s, transform 0.22s var(--ease);
}
.proof-card:hover {
  border-color: rgba(25,103,255,.3);
  transform: translateY(-3px);
}
.proof-header {
  display: flex;
  align-items: center;
  gap: 14px;
}
.proof-avatar {
  width: 52px; height: 52px;
  border-radius: 50%;
  display: flex;
  align-items: center;
  justify-content: center;
  font-family: var(--font-d);
  font-size: 20px;
  font-weight: 700;
  color: #fff;
  flex-shrink: 0;
}
.pa1 { background: linear-gradient(135deg,#1967ff,#3f82ff); }
.pa2 { background: linear-gradient(135deg,#ff6719,#ff9a5c); }
.pa3 { background: linear-gradient(135deg,#7c3aed,#a78bfa); }

.proof-name {
  font-family: var(--font-d);
  font-size: 16px;
  font-weight: 700;
  color: var(--ink1);
  letter-spacing: -0.01em;
}
.proof-company { font-size: 13px; color: var(--ink3); margin-top: 2px; }
.proof-metric {
  display: flex;
  align-items: baseline;
  gap: 6px;
}
.proof-num {
  font-family: var(--font-d);
  font-size: 36px;
  font-weight: 700;
  color: var(--accent);
  letter-spacing: -0.03em;
  line-height: 1;
}
.proof-unit { font-size: 14px; color: var(--ink3); }
\end{lstlisting}

The proof card is a vertical flex column with a 20-pixel gap, so its two halves
(the header and the metric) space themselves without margins.

\code{align-items: baseline} on line 594 is the detail worth learning. The metric
row holds a 36-pixel number beside 14-pixel text. Centring them would look wrong,
because the eye reads the two as sitting on a shared line of writing.
\code{baseline} aligns them on that invisible line the letters sit on, which is
what a typesetter would do.

The three avatar gradients (lines 579 to 581) repeat the same literal-colour issue
as the hero avatars, and \code{pa3} introduces two purples that exist nowhere in
the token block at all.

\paragraph{Lines 606 to 654: the services grid.}

\begin{lstlisting}[firstnumber=605,caption={\file{css/site.css}, lines 605--654}]

/* ─── SERVICES ─── */
.services-section { padding: 0 0 120px; }
.services-grid {
  display: grid;
  grid-template-columns: repeat(3, 1fr);
  gap: 2px;
  background: var(--border);
  border-radius: var(--r);
  overflow: hidden;
}
.service-item {
  background: var(--s0);
  padding: 40px 32px;
  display: flex;
  flex-direction: column;
  gap: 16px;
  transition: background 0.22s;
}
.service-item:hover { background: var(--s1); }
.service-icon {
  width: 52px; height: 52px;
  border-radius: var(--rsm);
  background: rgba(25,103,255,.1);
  border: 1px solid rgba(25,103,255,.2);
  display: flex;
  align-items: center;
  justify-content: center;
  transition: background 0.22s, border-color 0.22s;
}
.service-icon svg { width: 22px; height: 22px; color: var(--accentH); }
.service-item:hover .service-icon {
  background: rgba(25,103,255,.18);
  border-color: rgba(25,103,255,.45);
}
.service-name {
  font-family: var(--font-d);
  font-size: 17px;
  font-weight: 700;
  color: var(--ink1);
  letter-spacing: -0.01em;
}
.service-tag {
  display: inline-block;
  font-size: 11px;
  font-weight: 600;
  color: var(--ink3);
  text-transform: uppercase;
  letter-spacing: 0.07em;
}
\end{lstlisting}

Lines 608 to 615 use the technique mentioned in section 7.1 and it deserves its own
explanation, because it is genuinely clever and genuinely fragile.

\begin{keyidea}{Hairlines made of gaps}
The grid is given a 2-pixel gap and a background of the border colour. Each cell is
then given the page background. Because the cells cover the grid's background
everywhere except in the gaps, the gaps show through as 2-pixel lines. The result
is a perfect internal grid of hairlines with no border rules and no special case
for the last item in a row, which is exactly what the metrics band above needed
(\code{:last-child \{ border-right: none \}}).

The fragility is that it depends on the cell background matching the page. Put this
grid on a different surface without changing \code{.service-item}, and the cells
become visible rectangles of the wrong colour.
\end{keyidea}

\code{border-radius} plus \code{overflow: hidden} on lines 613 and 614 round the
whole grid and clip the corner cells to match.

Line 647 sets \code{display: inline-block} on the tag. Since the tag is inside a
flex column, it would already behave as a block, so this has no visible effect: it
is a harmless leftover.

\paragraph{Lines 656 to 725: the booking panel and the footer.}

\begin{lstlisting}[firstnumber=655,caption={\file{css/site.css}, lines 655--725}]

/* ─── BOOK A CALL ─── */
.book-section { padding: 0 0 140px; }
.book-wrap {
  background: var(--s1);
  border: 1px solid var(--border);
  border-radius: 24px;
  padding: 72px 48px;
  text-align: center;
  position: relative;
  overflow: hidden;
}
.book-wrap::before {
  content: '';
  position: absolute;
  top: -100px; left: 50%;
  transform: translateX(-50%);
  width: 600px; height: 300px;
  background: radial-gradient(ellipse, rgba(25,103,255,.12) 0%, transparent 70%);
  pointer-events: none;
}
.book-eyebrow {
  font-size: 12px;
  font-weight: 700;
  text-transform: uppercase;
  letter-spacing: 0.1em;
  color: var(--accent);
  margin-bottom: 16px;
}
.book-title {
  font-family: var(--font-d);
  font-size: clamp(28px, 4vw, 48px);
  font-weight: 700;
  letter-spacing: -0.025em;
  color: var(--ink1);
  margin-bottom: 12px;
}
.book-sub {
  font-size: 16px;
  color: var(--ink3);
  margin-bottom: 36px;
}
.book-cta-row {
  display: flex;
  align-items: center;
  justify-content: center;
  gap: 14px;
  flex-wrap: wrap;
}
.book-detail {
  display: flex;
  align-items: center;
  justify-content: center;
  gap: 6px;
  font-size: 13px;
  color: var(--ink3);
  margin-top: 20px;
}
.book-detail svg { width: 15px; height: 15px; }
.book-detail span { color: var(--ink2); }

/* ─── FOOTER ─── */
footer {
  border-top: 1px solid var(--borderS);
  padding: 32px 24px;
  text-align: center;
  font-size: 13px;
  color: var(--ink3);
}
footer a { color: var(--ink3); transition: color 0.15s; }
footer a:hover { color: var(--ink1); }
\end{lstlisting}

Lines 667 to 675 add a blue glow above the booking panel using an \code{::before}
pseudo-element: a 600 by 300 ellipse of brand blue at 12\% opacity, positioned 100
pixels above the panel's top edge and centred with the standard
\code{left: 50\%} plus \code{translateX(-50\%)} idiom. Because the panel has
\code{overflow: hidden}, only the bottom part of the glow is visible, which is what
makes it read as light spilling in from above rather than as a floating oval.

\code{pointer-events: none} on line 674 appears for the second time in the file, for
the same reason as before: this decorative layer sits over the panel and must not
intercept the click aimed at the booking button.

The footer is deliberately unremarkable: one hairline, centred dim text, and links
that brighten on hover over 0.15 seconds.

\paragraph{Lines 727 to 768: the responsive rules.}

\begin{lstlisting}[firstnumber=726,caption={\file{css/site.css}, lines 726--768}]

/* ─── RESPONSIVE ─── */
@media (max-width: 900px) {
  .metrics-grid { grid-template-columns: repeat(2, 1fr); }
  .metric-item { border-right: none; border-bottom: 1px solid var(--borderS); }
  .metric-item:nth-child(odd) { border-right: 1px solid var(--borderS); }
  .metric-item:nth-child(3),
  .metric-item:nth-child(4) { border-bottom: none; }

  .pipeline { flex-direction: column; }
  .pipeline-connector { width: 2px; height: 40px; }
  .pipeline-connector::after { width: 100% !important; height: 0; }
  .pipeline.visible .pipeline-connector::after { height: 100%; width: 100% !important; }
  .pipeline-lbl { white-space: normal; }

  .portfolio-grid { grid-template-columns: repeat(2, 1fr); }
  .proof-grid     { grid-template-columns: 1fr; }
  .services-grid  { grid-template-columns: repeat(2, 1fr); }

  .tile { display: none; }
  .tile:nth-child(1),
  .tile:nth-child(2),
  .tile:nth-child(3) { display: block; }
  .tile:nth-child(1) { top: 10%; left: 2%; width: 140px; }
  .tile:nth-child(2) { top: 15%; right: 2%; width: 140px; }
  .tile:nth-child(3) { display: none; }
}

@media (max-width: 600px) {
  .portfolio-grid { grid-template-columns: 1fr; }
  .services-grid  { grid-template-columns: 1fr; }
  .book-wrap { padding: 48px 24px; }
  .hero-h1 { font-size: 38px; }
  .tile { display: none; }
}

@media (prefers-reduced-motion: reduce) {
  *, *::before, *::after {
    animation-duration: .001ms !important;
    transition-duration: .001ms !important;
  }
  .reveal, .stagger > * { opacity: 1; transform: none; }
}
\end{lstlisting}

\begin{jargon}{Media query}
A block of CSS that only applies when a condition about the display is true.
\code{@media (max-width: 900px)} means ``apply everything in here only when the
window is 900 pixels wide or narrower''. Because it comes later in the file, its
rules override the equivalents above.
\end{jargon}

The 900-pixel block does five things. The metrics band folds from four columns to
two, and its dividing lines are rebuilt for the new shape: the right border is
removed from everything, a bottom border added to everything, the right border put
back on odd-numbered items so there is still a line down the middle, and the bottom
border removed from items 3 and 4 so the block does not end in a stray line. That
is four rules to reposition two lines, and it is the price of the border technique
that the services grid avoided.

The pipeline turns vertical (line 735) and its connectors become 2 pixels wide and
40 tall (line 736). Lines 737 and 738 then rebuild the fill animation to run
downwards instead of across, growing height rather than width, and use
\code{!important} to force the width value past the more specific rules further up
the file.

\begin{jargon}{!important}
A marker that forces a declaration to win against the normal precedence rules. It
is generally considered a last resort, because it breaks the predictability that
makes CSS manageable. It appears exactly twice in this file, both on line 737 and
738, both to override the horizontal fill animation.
\end{jargon}

The three content grids each drop a column, and the hero tiles are reduced from
eight to two.

\begin{honest}{Three lines of dead responsive code}
Line 745 hides every tile. Lines 746 to 748 then show tiles 1, 2 and 3. Line 751
hides tile 3 again. Because line 751 comes later and has identical specificity, it
wins, so tile 3 is hidden and the \code{display: block} it was given three lines
earlier does nothing at all.

The author evidently decided three tiles were too many for a tablet and removed the
third by adding a rule instead of editing the existing one. It is harmless, and it
is exactly the sort of thing that makes a reader doubt their understanding of the
cascade. It should be a two-line rule listing tiles 1 and 2.
\end{honest}

The 600-pixel block hides the tiles entirely, collapses the last two grids to one
column, tightens the booking panel's padding, and fixes the headline at 38 pixels,
overriding the \code{clamp} whose floor was 42.

The final block, lines 762 to 768, is the accessibility one, and it is the best
thing in the file.

\begin{keyidea}{What \code{prefers-reduced-motion} is and why this block matters}
Operating systems let people ask for reduced motion, usually because animation
causes them nausea, dizziness or migraine. The browser passes that preference to
the page, and this block responds to it.

It does two things. It crushes every animation and transition on the page to a
thousandth of a millisecond, which is effectively instant rather than disabled, a
deliberate choice because setting them to zero can stop some browsers firing the
events that mark an animation as finished. Then, crucially, line 767 forces every
revealed and staggered element to full opacity, because those elements start
invisible and would otherwise never appear once their transitions were neutered.

That second line is what turns a well-meaning block into a correct one, and it is
the detail most implementations of this pattern get wrong.
\end{keyidea}

\begin{honest}{The reduced-motion block cannot reach the counting numbers}
This block governs CSS only. The four metric numbers are animated by JavaScript,
counting up over 2.2 seconds, and \file{main.js} never checks this preference. A
visitor who has asked their system for reduced motion still gets four numbers
whirring upward.

The fix is four lines in \file{main.js}, given in section 7.3. This is the single
clearest defect in the project: not because it is severe, but because the intent to
handle it correctly is demonstrably there in the CSS and simply was not carried
across to the other file.
\end{honest}

\subsubsection*{What it connects to}

\textbf{Named by:} \file{new-design-reference/index.html}, line 16, and nothing
else.

\textbf{Names:} no other file. It refers to no images, no fonts by file (it names
font \emph{families}, which the HTML is responsible for downloading), and imports
nothing.

\textbf{Coupled to \file{index.html}} in three ways that no tool checks: by class
name throughout, by child position for the hero tiles (lines 162 to 169) and the
pipeline (lines 400 to 426), and by the shared contract that JavaScript adds the
class \code{visible}.

\textbf{Coupled to \file{js/main.js}} through exactly two class names,
\code{scrolled} and \code{visible}, both of which the JavaScript adds and this file
defines.

\subsubsection*{Concepts introduced}

This file is where most of the page's technical vocabulary first appears: custom
properties, the cascade, flexbox, grid, pseudo-elements, pseudo-classes, keyframe
animations, transitions, media queries, \code{clamp}, \code{calc},
\code{aspect-ratio}, \code{backdrop-filter}, stacking with \code{z-index}, and
\code{prefers-reduced-motion}. Each is defined at its first use above and collected
in the glossary.

\subsubsection*{How it breaks}

\begin{itemize}
\item \textbf{Delete this file and the page still works}, in the sense that every
  word remains readable. It simply looks like an unstyled document, and the hero
  tiles appear as a plain stacked list.
\item \textbf{Rename a class here without renaming it in the HTML} and that
  element silently loses its styling. Nothing reports it.
\item \textbf{Add or remove a hero tile or a pipeline step} and the position rules
  stop lining up, silently.
\item \textbf{Remove \code{pointer-events: none} on line 149} and the hero tiles
  begin swallowing clicks meant for the buttons beneath them.
\item \textbf{Remove \code{position: relative} on line 137} and all eight tiles
  jump to positions measured against the whole window instead of the hero.
\item \textbf{A missing semicolon} makes the browser discard the rest of that
  declaration block. Unlike a program, CSS never reports an error: it drops what it
  cannot parse and carries on, which is why a stylesheet bug shows up as something
  looking wrong rather than as a message.
\end{itemize}

\subsubsection*{Honest findings for this file}

Two unused tokens (lines 14 and 24); one dead rule with no matching element (line
352); one rule permanently overridden by inline styles (line 209); three lines of
dead responsive code (lines 746 to 748 versus 751); a doubled flex layout in the
navigation (lines 100 to 102 and 116 to 118); smooth scrolling implemented twice,
here and in JavaScript; literal colours used where tokens exist (lines 319 to 322,
579 to 581); two different techniques for drawing dividing lines in one file; no
focus states on either button; and a reduced-motion block that is exemplary within
CSS and cannot reach the one animation that lives outside it.
